\documentclass[11pt, a4paper]{article}
% ── Shared preamble for all RTOS lab documents ────────────────────────────────
% Usage: % ── Shared preamble for all RTOS lab documents ────────────────────────────────
% Usage: % ── Shared preamble for all RTOS lab documents ────────────────────────────────
% Usage: \input{../preamble}  (from each labNN/ directory)
% ──────────────────────────────────────────────────────────────────────────────

% ── Packages ──────────────────────────────────────────────────────────────────
\usepackage[a4paper, top=2.5cm, bottom=2.5cm, left=2.5cm, right=2.5cm, headheight=28pt]{geometry}
\usepackage[utf8]{inputenc}
\usepackage[T1]{fontenc}
\usepackage{lmodern}
\usepackage{booktabs}
\usepackage{array}
\usepackage{tabularx}
\usepackage{longtable}
\usepackage{multirow}
\usepackage{colortbl}
\usepackage{xcolor}
\usepackage{titlesec}
\usepackage{enumitem}
\usepackage{hyperref}
\usepackage{fancyhdr}
\usepackage{parskip}
\usepackage{listings}
\usepackage[most]{tcolorbox}
\usepackage{fontawesome5}
\usepackage{microtype}
\usepackage{graphicx}
\usepackage{amsmath}

% ── Colors (KMUTNB brand) ──────────────────────────────────────────────────────
\definecolor{kmutnbBlue}{RGB}{0,56,116}
\definecolor{kmutnbGold}{RGB}{186,146,36}
\definecolor{lightGray}{RGB}{245,245,245}
\definecolor{midGray}{RGB}{200,200,200}
\definecolor{codeGray}{RGB}{248,248,248}
\definecolor{codeFrame}{RGB}{220,220,220}
\definecolor{tipteal}{RGB}{0,105,92}
\definecolor{warnAmber}{RGB}{121,85,0}
\definecolor{checkGreen}{RGB}{27,94,32}

% ── Section formatting ─────────────────────────────────────────────────────────
\titleformat{\section}
  {\large\bfseries\color{kmutnbBlue}}
  {\thesection.}{0.5em}{}
  [\color{kmutnbGold}\titlerule]

\titleformat{\subsection}
  {\normalsize\bfseries\color{kmutnbBlue}}
  {\thesubsection.}{0.5em}{}

\titleformat{\subsubsection}
  {\small\bfseries\color{kmutnbBlue}}
  {\thesubsubsection.}{0.5em}{}

\titlespacing*{\section}{0pt}{1.4ex plus .2ex}{0.6ex plus .1ex}
\titlespacing*{\subsection}{0pt}{1.0ex plus .2ex}{0.4ex plus .1ex}

% ── Listings (C and bash) ──────────────────────────────────────────────────────
\lstdefinestyle{cstyle}{
  language        = C,
  basicstyle      = \ttfamily\small,
  keywordstyle    = \bfseries\color{kmutnbBlue},
  commentstyle    = \itshape\color{gray},
  stringstyle     = \color{tipteal},
  numberstyle     = \tiny\color{midGray},
  numbers         = left,
  numbersep       = 6pt,
  stepnumber      = 1,
  backgroundcolor = \color{codeGray},
  frame           = single,
  framerule       = 0.4pt,
  rulecolor       = \color{codeFrame},
  breaklines      = true,
  showstringspaces= false,
  tabsize         = 4,
  xleftmargin     = 1.5em,
  xrightmargin    = 0.5em,
  aboveskip       = 0.8em,
  belowskip       = 0.8em,
}

\lstdefinestyle{bashstyle}{
  language        = bash,
  basicstyle      = \ttfamily\small,
  keywordstyle    = \color{kmutnbBlue},
  commentstyle    = \itshape\color{gray},
  backgroundcolor = \color{black!5},
  frame           = single,
  framerule       = 0.4pt,
  rulecolor       = \color{codeFrame},
  breaklines      = true,
  showstringspaces= false,
  xleftmargin     = 1.5em,
  xrightmargin    = 0.5em,
  aboveskip       = 0.6em,
  belowskip       = 0.6em,
}

\lstdefinestyle{textstyle}{
  basicstyle      = \ttfamily\footnotesize,
  backgroundcolor = \color{black!4},
  frame           = single,
  framerule       = 0.4pt,
  rulecolor       = \color{codeFrame},
  breaklines      = true,
  xleftmargin     = 1.5em,
  xrightmargin    = 0.5em,
  aboveskip       = 0.6em,
  belowskip       = 0.6em,
}

% ── Callout boxes ──────────────────────────────────────────────────────────────
\tcbuselibrary{skins, breakable}

\newtcolorbox{tipbox}[1][]{
  enhanced, breakable,
  colback=tipteal!8, colframe=tipteal, coltitle=white,
  fonttitle=\bfseries\small, title={\faLightbulb\quad Tip#1},
  left=6pt, right=6pt, top=4pt, bottom=4pt,
  boxrule=0.8pt, arc=3pt,
}

\newtcolorbox{warnbox}[1][]{
  enhanced, breakable,
  colback=warnAmber!8, colframe=kmutnbGold, coltitle=white,
  fonttitle=\bfseries\small, title={\faExclamationTriangle\quad Note#1},
  left=6pt, right=6pt, top=4pt, bottom=4pt,
  boxrule=0.8pt, arc=3pt,
}

\newtcolorbox{checkpointbox}[1][]{
  enhanced,
  colback=kmutnbBlue!6, colframe=kmutnbBlue, coltitle=white,
  fonttitle=\bfseries\small, title={\faCheckCircle\quad Checkpoint#1},
  left=6pt, right=6pt, top=4pt, bottom=4pt,
  boxrule=0.8pt, arc=3pt,
}

\newtcolorbox{questionbox}[1][]{
  enhanced, breakable,
  colback=kmutnbGold!10, colframe=kmutnbGold, coltitle=white,
  fonttitle=\bfseries\small, title={\faQuestion\quad Question#1},
  left=6pt, right=6pt, top=4pt, bottom=4pt,
  boxrule=0.8pt, arc=3pt,
}

% ── Checkbox list ──────────────────────────────────────────────────────────────
\newlist{checklist}{itemize}{1}
\setlist[checklist]{label=\faSquare[regular], leftmargin=1.5em, noitemsep}

% ── Misc helpers ───────────────────────────────────────────────────────────────
\newcommand{\file}[1]{\texttt{\small #1}}
\newcommand{\api}[1]{\texttt{\small #1}}

% ── Title block macro ─────────────────────────────────────────────────────────
% Usage: \labtitle{03}{Aperiodic Servers \& Producer--Consumer}{2}{CLO 1 \& CLO 2}
\newcommand{\labtitle}[4]{%
  \begin{center}
    {\color{kmutnbBlue}\rule{\linewidth}{2pt}}\\[6pt]
    {\LARGE\bfseries\color{kmutnbBlue} Laboratory #1}\\[4pt]
    {\large\bfseries #2}\\[2pt]
    {\normalsize Real-Time Operating Systems \enspace\textbullet\enspace
     M.Eng.\ ECE \enspace\textbullet\enspace KMUTNB}\\[8pt]
    {\color{kmutnbGold}\rule{\linewidth}{1pt}}
  \end{center}
  \vspace{0.3cm}
  \begin{tabular}{@{}p{3.8cm} p{10cm}@{}}
    \textbf{Platform}       & NXP FRDM-MCXN236 (ARM Cortex-M33 @ 150\,MHz) \\[2pt]
    \textbf{RTOS}           & FreeRTOS v10.6 \\[2pt]
    \textbf{Toolchain}      & ARM GNU Toolchain (\texttt{arm-none-eabi-gcc}) + Make \\[2pt]
    \textbf{Estimated time} & #3 hours \\[2pt]
    \textbf{CLO mapping}    & #4 \\
  \end{tabular}
  \vspace{0.4cm}
}
  (from each labNN/ directory)
% ──────────────────────────────────────────────────────────────────────────────

% ── Packages ──────────────────────────────────────────────────────────────────
\usepackage[a4paper, top=2.5cm, bottom=2.5cm, left=2.5cm, right=2.5cm, headheight=28pt]{geometry}
\usepackage[utf8]{inputenc}
\usepackage[T1]{fontenc}
\usepackage{lmodern}
\usepackage{booktabs}
\usepackage{array}
\usepackage{tabularx}
\usepackage{longtable}
\usepackage{multirow}
\usepackage{colortbl}
\usepackage{xcolor}
\usepackage{titlesec}
\usepackage{enumitem}
\usepackage{hyperref}
\usepackage{fancyhdr}
\usepackage{parskip}
\usepackage{listings}
\usepackage[most]{tcolorbox}
\usepackage{fontawesome5}
\usepackage{microtype}
\usepackage{graphicx}
\usepackage{amsmath}

% ── Colors (KMUTNB brand) ──────────────────────────────────────────────────────
\definecolor{kmutnbBlue}{RGB}{0,56,116}
\definecolor{kmutnbGold}{RGB}{186,146,36}
\definecolor{lightGray}{RGB}{245,245,245}
\definecolor{midGray}{RGB}{200,200,200}
\definecolor{codeGray}{RGB}{248,248,248}
\definecolor{codeFrame}{RGB}{220,220,220}
\definecolor{tipteal}{RGB}{0,105,92}
\definecolor{warnAmber}{RGB}{121,85,0}
\definecolor{checkGreen}{RGB}{27,94,32}

% ── Section formatting ─────────────────────────────────────────────────────────
\titleformat{\section}
  {\large\bfseries\color{kmutnbBlue}}
  {\thesection.}{0.5em}{}
  [\color{kmutnbGold}\titlerule]

\titleformat{\subsection}
  {\normalsize\bfseries\color{kmutnbBlue}}
  {\thesubsection.}{0.5em}{}

\titleformat{\subsubsection}
  {\small\bfseries\color{kmutnbBlue}}
  {\thesubsubsection.}{0.5em}{}

\titlespacing*{\section}{0pt}{1.4ex plus .2ex}{0.6ex plus .1ex}
\titlespacing*{\subsection}{0pt}{1.0ex plus .2ex}{0.4ex plus .1ex}

% ── Listings (C and bash) ──────────────────────────────────────────────────────
\lstdefinestyle{cstyle}{
  language        = C,
  basicstyle      = \ttfamily\small,
  keywordstyle    = \bfseries\color{kmutnbBlue},
  commentstyle    = \itshape\color{gray},
  stringstyle     = \color{tipteal},
  numberstyle     = \tiny\color{midGray},
  numbers         = left,
  numbersep       = 6pt,
  stepnumber      = 1,
  backgroundcolor = \color{codeGray},
  frame           = single,
  framerule       = 0.4pt,
  rulecolor       = \color{codeFrame},
  breaklines      = true,
  showstringspaces= false,
  tabsize         = 4,
  xleftmargin     = 1.5em,
  xrightmargin    = 0.5em,
  aboveskip       = 0.8em,
  belowskip       = 0.8em,
}

\lstdefinestyle{bashstyle}{
  language        = bash,
  basicstyle      = \ttfamily\small,
  keywordstyle    = \color{kmutnbBlue},
  commentstyle    = \itshape\color{gray},
  backgroundcolor = \color{black!5},
  frame           = single,
  framerule       = 0.4pt,
  rulecolor       = \color{codeFrame},
  breaklines      = true,
  showstringspaces= false,
  xleftmargin     = 1.5em,
  xrightmargin    = 0.5em,
  aboveskip       = 0.6em,
  belowskip       = 0.6em,
}

\lstdefinestyle{textstyle}{
  basicstyle      = \ttfamily\footnotesize,
  backgroundcolor = \color{black!4},
  frame           = single,
  framerule       = 0.4pt,
  rulecolor       = \color{codeFrame},
  breaklines      = true,
  xleftmargin     = 1.5em,
  xrightmargin    = 0.5em,
  aboveskip       = 0.6em,
  belowskip       = 0.6em,
}

% ── Callout boxes ──────────────────────────────────────────────────────────────
\tcbuselibrary{skins, breakable}

\newtcolorbox{tipbox}[1][]{
  enhanced, breakable,
  colback=tipteal!8, colframe=tipteal, coltitle=white,
  fonttitle=\bfseries\small, title={\faLightbulb\quad Tip#1},
  left=6pt, right=6pt, top=4pt, bottom=4pt,
  boxrule=0.8pt, arc=3pt,
}

\newtcolorbox{warnbox}[1][]{
  enhanced, breakable,
  colback=warnAmber!8, colframe=kmutnbGold, coltitle=white,
  fonttitle=\bfseries\small, title={\faExclamationTriangle\quad Note#1},
  left=6pt, right=6pt, top=4pt, bottom=4pt,
  boxrule=0.8pt, arc=3pt,
}

\newtcolorbox{checkpointbox}[1][]{
  enhanced,
  colback=kmutnbBlue!6, colframe=kmutnbBlue, coltitle=white,
  fonttitle=\bfseries\small, title={\faCheckCircle\quad Checkpoint#1},
  left=6pt, right=6pt, top=4pt, bottom=4pt,
  boxrule=0.8pt, arc=3pt,
}

\newtcolorbox{questionbox}[1][]{
  enhanced, breakable,
  colback=kmutnbGold!10, colframe=kmutnbGold, coltitle=white,
  fonttitle=\bfseries\small, title={\faQuestion\quad Question#1},
  left=6pt, right=6pt, top=4pt, bottom=4pt,
  boxrule=0.8pt, arc=3pt,
}

% ── Checkbox list ──────────────────────────────────────────────────────────────
\newlist{checklist}{itemize}{1}
\setlist[checklist]{label=\faSquare[regular], leftmargin=1.5em, noitemsep}

% ── Misc helpers ───────────────────────────────────────────────────────────────
\newcommand{\file}[1]{\texttt{\small #1}}
\newcommand{\api}[1]{\texttt{\small #1}}

% ── Title block macro ─────────────────────────────────────────────────────────
% Usage: \labtitle{03}{Aperiodic Servers \& Producer--Consumer}{2}{CLO 1 \& CLO 2}
\newcommand{\labtitle}[4]{%
  \begin{center}
    {\color{kmutnbBlue}\rule{\linewidth}{2pt}}\\[6pt]
    {\LARGE\bfseries\color{kmutnbBlue} Laboratory #1}\\[4pt]
    {\large\bfseries #2}\\[2pt]
    {\normalsize Real-Time Operating Systems \enspace\textbullet\enspace
     M.Eng.\ ECE \enspace\textbullet\enspace KMUTNB}\\[8pt]
    {\color{kmutnbGold}\rule{\linewidth}{1pt}}
  \end{center}
  \vspace{0.3cm}
  \begin{tabular}{@{}p{3.8cm} p{10cm}@{}}
    \textbf{Platform}       & NXP FRDM-MCXN236 (ARM Cortex-M33 @ 150\,MHz) \\[2pt]
    \textbf{RTOS}           & FreeRTOS v10.6 \\[2pt]
    \textbf{Toolchain}      & ARM GNU Toolchain (\texttt{arm-none-eabi-gcc}) + Make \\[2pt]
    \textbf{Estimated time} & #3 hours \\[2pt]
    \textbf{CLO mapping}    & #4 \\
  \end{tabular}
  \vspace{0.4cm}
}
  (from each labNN/ directory)
% ──────────────────────────────────────────────────────────────────────────────

% ── Packages ──────────────────────────────────────────────────────────────────
\usepackage[a4paper, top=2.5cm, bottom=2.5cm, left=2.5cm, right=2.5cm, headheight=28pt]{geometry}
\usepackage[utf8]{inputenc}
\usepackage[T1]{fontenc}
\usepackage{lmodern}
\usepackage{booktabs}
\usepackage{array}
\usepackage{tabularx}
\usepackage{longtable}
\usepackage{multirow}
\usepackage{colortbl}
\usepackage{xcolor}
\usepackage{titlesec}
\usepackage{enumitem}
\usepackage{hyperref}
\usepackage{fancyhdr}
\usepackage{parskip}
\usepackage{listings}
\usepackage[most]{tcolorbox}
\usepackage{fontawesome5}
\usepackage{microtype}
\usepackage{graphicx}
\usepackage{amsmath}

% ── Colors (KMUTNB brand) ──────────────────────────────────────────────────────
\definecolor{kmutnbBlue}{RGB}{0,56,116}
\definecolor{kmutnbGold}{RGB}{186,146,36}
\definecolor{lightGray}{RGB}{245,245,245}
\definecolor{midGray}{RGB}{200,200,200}
\definecolor{codeGray}{RGB}{248,248,248}
\definecolor{codeFrame}{RGB}{220,220,220}
\definecolor{tipteal}{RGB}{0,105,92}
\definecolor{warnAmber}{RGB}{121,85,0}
\definecolor{checkGreen}{RGB}{27,94,32}

% ── Section formatting ─────────────────────────────────────────────────────────
\titleformat{\section}
  {\large\bfseries\color{kmutnbBlue}}
  {\thesection.}{0.5em}{}
  [\color{kmutnbGold}\titlerule]

\titleformat{\subsection}
  {\normalsize\bfseries\color{kmutnbBlue}}
  {\thesubsection.}{0.5em}{}

\titleformat{\subsubsection}
  {\small\bfseries\color{kmutnbBlue}}
  {\thesubsubsection.}{0.5em}{}

\titlespacing*{\section}{0pt}{1.4ex plus .2ex}{0.6ex plus .1ex}
\titlespacing*{\subsection}{0pt}{1.0ex plus .2ex}{0.4ex plus .1ex}

% ── Listings (C and bash) ──────────────────────────────────────────────────────
\lstdefinestyle{cstyle}{
  language        = C,
  basicstyle      = \ttfamily\small,
  keywordstyle    = \bfseries\color{kmutnbBlue},
  commentstyle    = \itshape\color{gray},
  stringstyle     = \color{tipteal},
  numberstyle     = \tiny\color{midGray},
  numbers         = left,
  numbersep       = 6pt,
  stepnumber      = 1,
  backgroundcolor = \color{codeGray},
  frame           = single,
  framerule       = 0.4pt,
  rulecolor       = \color{codeFrame},
  breaklines      = true,
  showstringspaces= false,
  tabsize         = 4,
  xleftmargin     = 1.5em,
  xrightmargin    = 0.5em,
  aboveskip       = 0.8em,
  belowskip       = 0.8em,
}

\lstdefinestyle{bashstyle}{
  language        = bash,
  basicstyle      = \ttfamily\small,
  keywordstyle    = \color{kmutnbBlue},
  commentstyle    = \itshape\color{gray},
  backgroundcolor = \color{black!5},
  frame           = single,
  framerule       = 0.4pt,
  rulecolor       = \color{codeFrame},
  breaklines      = true,
  showstringspaces= false,
  xleftmargin     = 1.5em,
  xrightmargin    = 0.5em,
  aboveskip       = 0.6em,
  belowskip       = 0.6em,
}

\lstdefinestyle{textstyle}{
  basicstyle      = \ttfamily\footnotesize,
  backgroundcolor = \color{black!4},
  frame           = single,
  framerule       = 0.4pt,
  rulecolor       = \color{codeFrame},
  breaklines      = true,
  xleftmargin     = 1.5em,
  xrightmargin    = 0.5em,
  aboveskip       = 0.6em,
  belowskip       = 0.6em,
}

% ── Callout boxes ──────────────────────────────────────────────────────────────
\tcbuselibrary{skins, breakable}

\newtcolorbox{tipbox}[1][]{
  enhanced, breakable,
  colback=tipteal!8, colframe=tipteal, coltitle=white,
  fonttitle=\bfseries\small, title={\faLightbulb\quad Tip#1},
  left=6pt, right=6pt, top=4pt, bottom=4pt,
  boxrule=0.8pt, arc=3pt,
}

\newtcolorbox{warnbox}[1][]{
  enhanced, breakable,
  colback=warnAmber!8, colframe=kmutnbGold, coltitle=white,
  fonttitle=\bfseries\small, title={\faExclamationTriangle\quad Note#1},
  left=6pt, right=6pt, top=4pt, bottom=4pt,
  boxrule=0.8pt, arc=3pt,
}

\newtcolorbox{checkpointbox}[1][]{
  enhanced,
  colback=kmutnbBlue!6, colframe=kmutnbBlue, coltitle=white,
  fonttitle=\bfseries\small, title={\faCheckCircle\quad Checkpoint#1},
  left=6pt, right=6pt, top=4pt, bottom=4pt,
  boxrule=0.8pt, arc=3pt,
}

\newtcolorbox{questionbox}[1][]{
  enhanced, breakable,
  colback=kmutnbGold!10, colframe=kmutnbGold, coltitle=white,
  fonttitle=\bfseries\small, title={\faQuestion\quad Question#1},
  left=6pt, right=6pt, top=4pt, bottom=4pt,
  boxrule=0.8pt, arc=3pt,
}

% ── Checkbox list ──────────────────────────────────────────────────────────────
\newlist{checklist}{itemize}{1}
\setlist[checklist]{label=\faSquare[regular], leftmargin=1.5em, noitemsep}

% ── Misc helpers ───────────────────────────────────────────────────────────────
\newcommand{\file}[1]{\texttt{\small #1}}
\newcommand{\api}[1]{\texttt{\small #1}}

% ── Title block macro ─────────────────────────────────────────────────────────
% Usage: \labtitle{03}{Aperiodic Servers \& Producer--Consumer}{2}{CLO 1 \& CLO 2}
\newcommand{\labtitle}[4]{%
  \begin{center}
    {\color{kmutnbBlue}\rule{\linewidth}{2pt}}\\[6pt]
    {\LARGE\bfseries\color{kmutnbBlue} Laboratory #1}\\[4pt]
    {\large\bfseries #2}\\[2pt]
    {\normalsize Real-Time Operating Systems \enspace\textbullet\enspace
     M.Eng.\ ECE \enspace\textbullet\enspace KMUTNB}\\[8pt]
    {\color{kmutnbGold}\rule{\linewidth}{1pt}}
  \end{center}
  \vspace{0.3cm}
  \begin{tabular}{@{}p{3.8cm} p{10cm}@{}}
    \textbf{Platform}       & NXP FRDM-MCXN236 (ARM Cortex-M33 @ 150\,MHz) \\[2pt]
    \textbf{RTOS}           & FreeRTOS v10.6 \\[2pt]
    \textbf{Toolchain}      & ARM GNU Toolchain (\texttt{arm-none-eabi-gcc}) + Make \\[2pt]
    \textbf{Estimated time} & #3 hours \\[2pt]
    \textbf{CLO mapping}    & #4 \\
  \end{tabular}
  \vspace{0.4cm}
}


% ── Header / Footer ────────────────────────────────────────────────────────────
\pagestyle{fancy}
\fancyhf{}
\fancyhead[L]{\small\color{kmutnbBlue}\textbf{KMUTNB -- Faculty of Engineering}}
\fancyhead[R]{\small\color{kmutnbBlue}Dept.\ of Electrical and Computer Engineering}
\fancyfoot[L]{\small\color{kmutnbBlue}Real-Time Operating Systems -- M.Eng.\ ECE}
\fancyfoot[C]{\small\thepage}
\fancyfoot[R]{\small\color{kmutnbBlue}Lab 09}
\renewcommand{\headrulewidth}{0.4pt}
\renewcommand{\headrule}{\hbox to\headwidth{%
  \color{kmutnbGold}\leaders\hrule height\headrulewidth\hfill}}
\renewcommand{\footrulewidth}{0.4pt}
\renewcommand{\footrule}{\hbox to\headwidth{%
  \color{midGray}\leaders\hrule height\footrulewidth\hfill}}

\hypersetup{
  colorlinks = true,
  linkcolor  = kmutnbBlue,
  urlcolor   = kmutnbBlue,
  citecolor  = kmutnbBlue,
  pdftitle   = {Lab 09 -- Dual-Core AMP: Messaging Unit and Shared SRAM},
  pdfauthor  = {KMUTNB RTOS Course},
}

% ══════════════════════════════════════════════════════════════════════════════
\begin{document}

\labtitle{09}{Dual-Core AMP: Messaging Unit \& Shared SRAM}{4--5}{CLO 5 \& CLO 6}

% ── Objectives ────────────────────────────────────────────────────────────────
\section{Objectives}

By completing this lab you will be able to:

\begin{enumerate}[noitemsep, leftmargin=*]
  \item \textbf{Boot} both CM33 cores of the MCXN236 from a single project using SYSCON
        boot address registers.
  \item \textbf{Exchange} 32-bit messages between cores using the Messaging Unit (MU)
        mailboxes.
  \item \textbf{Implement} a shared SRAM ring buffer with correct DMB memory barriers.
  \item \textbf{Measure} inter-core IPC latency (MU send $\to$ ISR entry) using DWT
        timestamps on both cores.
  \item \textbf{Compare} MU-direct vs.\ ring-buffer throughput for high-rate data
        transfer.
\end{enumerate}

% ── Prerequisites ─────────────────────────────────────────────────────────────
\section{Prerequisites}

\begin{checklist}
  \item Lab 06 complete (DWT cycle counter working).
  \item ARM GNU Toolchain and \texttt{pyocd} working.
  \item MCUXpresso SDK for MCXN236 (for \file{fsl\_mu.h} and multicore startup code).
  \item Understanding of memory barriers (\api{\_\_DMB()}, \api{\_\_DSB()}).
\end{checklist}

% ── Background ────────────────────────────────────────────────────────────────
\section{Background}

\subsection{MCXN236 AMP Configuration}

The MCXN236 has two independent Cortex-M33 cores (CM33\_0 and CM33\_1). In AMP mode:

\begin{itemize}[noitemsep, leftmargin=*]
  \item Each core runs its own independent firmware image with its own FreeRTOS instance.
  \item CM33\_0 boots first (primary core). It writes the CM33\_1 boot address and
        releases the reset.
  \item Communication uses the \textbf{Messaging Unit (MU)} --- four 32-bit TX/RX
        registers per direction, with hardware flow control.
  \item Shared data lives in \textbf{SRAM4} (0x2008\_0000, 128\,KB) --- accessible by
        both cores.
\end{itemize}

\subsection{Memory Barriers}

The Cortex-M33 has a weakly ordered memory model. Without barriers, the CPU may
reorder stores before the MU send, causing the consumer to read stale data:

\begin{lstlisting}[style=cstyle]
/* Producer -- on CM33_0 */
shared_buf[idx] = data;   /* store */
__DMB();                   /* all preceding stores visible before MU send */
MU_SendMsg(MUA, 0, idx);   /* notify consumer */

/* Consumer ISR -- on CM33_1 */
uint32_t idx = MU_ReceiveMsg(MUB, 0);
__DMB();                   /* MU read complete before we read shared_buf */
process(shared_buf[idx]);
\end{lstlisting}

% ── Project Structure ─────────────────────────────────────────────────────────
\section{Project Structure}

\begin{lstlisting}[style=textstyle]
lab09_dual_core/
  cm33_0/                          <- Primary core firmware
    Makefile
    linker_cm33_0.ld               <- SRAM0+1 private, SRAM4 shared
    src/
      main_cm33_0.c                <- boots CM33_1, creates tasks
      adc_producer.c/.h            <- ADC sampling task (1 kHz)
      shared_mem.h                 <- shared ring buffer declaration
      FreeRTOSConfig.h
      uart.h / uart.c
  cm33_1/                          <- Secondary core firmware
    Makefile
    linker_cm33_1.ld               <- SRAM2+3 private, SRAM4 shared
    src/
      main_cm33_1.c                <- MU init, FreeRTOS start
      fir_consumer.c/.h            <- FIR filter task
      shared_mem.h                 <- same shared ring buffer header
      FreeRTOSConfig.h
      uart.h / uart.c
  flash_dual.sh                    <- flashes both images
\end{lstlisting}

% ── Part A ─────────────────────────────────────────────────────────────────────
\section{Part A --- Boot Both Cores}

\subsection{Step 1 --- Release CM33\_1 from CM33\_0}

In \file{cm33\_0/src/main\_cm33\_0.c}, after hardware init and before starting the
scheduler:

\begin{lstlisting}[style=cstyle]
#include "fsl_syscon.h"

void Boot_CM33_1(void)
{
    /* Write CM33_1's vector table address (= start of its flash region) */
    SYSCON->CPBOOT   = 0x00080000U;
    SYSCON->CPUCTRL |= SYSCON_CPUCTRL_CPU1RSTEN_MASK;   /* assert reset */
    SYSCON->CPUCTRL &= ~SYSCON_CPUCTRL_CPU1RSTEN_MASK;  /* release reset */
    uart_printf("[CM0] CM33_1 released at 0x00080000\r\n");
}
\end{lstlisting}

\subsection{Step 2 --- CM33\_1 Startup}

CM33\_1's \file{main\_cm33\_1.c} initialises its own hardware and announces its
presence:

\begin{lstlisting}[style=cstyle]
int main(void)
{
    BOARD_InitHardware_CM33_1();
    uart_printf("[CM1] Core 1 running at %lu Hz\r\n",
                CLOCK_GetFreq(kCLOCK_CoreSysClk));
    /* ... task creation ... */
    vTaskStartScheduler();
}
\end{lstlisting}

\begin{checkpointbox}[~A]
  UART showing both \texttt{[CM0]} and \texttt{[CM1]} startup messages.
\end{checkpointbox}

\begin{questionbox}[~A1]
  CM33\_0 releases CM33\_1 by writing \texttt{CPBOOT} and toggling
  \texttt{CPU1RSTEN}. What would happen if CM33\_0 wrote \texttt{CPBOOT} but never
  cleared \texttt{CPU1RSTEN}?
\end{questionbox}

% ── Part B ─────────────────────────────────────────────────────────────────────
\section{Part B --- MU Direct Messaging}

\subsection{Step 1 --- Send 32-Bit Messages}

On CM33\_0 (sender task):

\begin{lstlisting}[style=cstyle]
void vMuSendTask(void *pv)
{
    uint32_t count = 0;
    for (;;) {
        vTaskDelay(pdMS_TO_TICKS(100));
        MU_SendMsg(MUA, 0, count);
        uart_printf("[CM0] sent %lu\r\n", count);
        count++;
    }
}
\end{lstlisting}

On CM33\_1 (MU receive ISR):

\begin{lstlisting}[style=cstyle]
void MUA_IRQHandler(void)
{
    uint32_t msg = MU_ReceiveMsg(MUB, 0);
    MU_ClearStatusFlags(MUB, kMU_Rx0FullFlag);
    uart_printf("[CM1] received %lu\r\n", msg);
}

/* In main_cm33_1.c: enable MU RX interrupt */
MU_EnableInterrupts(MUB, kMU_Rx0FullInterruptEnable);
EnableIRQ(MUA_IRQn);
\end{lstlisting}

\begin{checkpointbox}[~B]
  UART showing matching send/receive sequence from both cores.
\end{checkpointbox}

\subsection{Step 2 --- Measure MU Latency with DWT}

On CM33\_0, record the DWT cycle count just before \api{MU\_SendMsg}. On CM33\_1,
record it in the first instruction of the ISR. Use a shared SRAM4 variable:

\begin{lstlisting}[style=cstyle]
/* shared_mem.h */
typedef struct {
    volatile uint32_t send_cycles;    /* written by CM33_0 */
    volatile uint32_t recv_cycles;    /* written by CM33_1 ISR */
} __attribute__((packed)) LatencyTest_t;

extern LatencyTest_t g_latency;   /* in .sram4 section */
\end{lstlisting}

\begin{lstlisting}[style=cstyle]
/* CM33_0 sender */
g_latency.send_cycles = DWT->CYCCNT;
__DMB();
MU_SendMsg(MUA, 0, 0xDEAD);

/* CM33_1 ISR -- first two lines */
g_latency.recv_cycles = DWT->CYCCNT;
__DMB();
\end{lstlisting}

After a few exchanges, CM33\_1 prints the latency:

\begin{lstlisting}[style=cstyle]
uint32_t latency = g_latency.recv_cycles - g_latency.send_cycles;
uart_printf("[CM1] MU latency = %lu cycles (%.1f ns @ 150 MHz)\r\n",
            latency, latency / 150.0f);
\end{lstlisting}

\begin{checkpointbox}[~B2]
  Terminal showing measured MU latency in cycles and nanoseconds.
\end{checkpointbox}

\begin{questionbox}[~B1]
  The MU TX register has hardware flow control --- \api{MU\_SendMsg} blocks until the
  remote core reads the register. For a high-frequency producer (1\,kHz), is this
  acceptable? How would you fix it without blocking the sender?
\end{questionbox}

% ── Part C ─────────────────────────────────────────────────────────────────────
\section{Part C --- Shared SRAM Ring Buffer}

\subsection{Step 1 --- Place Ring Buffer in SRAM4}

In \file{shared\_mem.h} (identical copy in both projects):

\begin{lstlisting}[style=cstyle]
#define RING_BUF_SIZE  256

typedef struct {
    volatile uint32_t write_idx;
    volatile uint32_t read_idx;
    volatile uint32_t missed;      /* overflow counter */
    int16_t           data[RING_BUF_SIZE];
} __attribute__((packed, aligned(4))) RingBuf_t;

extern RingBuf_t g_ring;   /* placed in .sram4 section */
\end{lstlisting}

In the linker script \file{cm33\_0/linker\_cm33\_0.ld}:

\begin{lstlisting}[style=textstyle]
MEMORY {
  flash0  (rx)  : ORIGIN = 0x00000000, LENGTH = 512K
  sram01  (rwx) : ORIGIN = 0x20000000, LENGTH = 256K   /* CM33_0 private */
  sram4   (rw)  : ORIGIN = 0x20080000, LENGTH = 128K   /* shared */
}

SECTIONS {
  .sram4 (NOLOAD) : {
    KEEP(*(.sram4*))
  } > sram4
}
\end{lstlisting}

\subsection{Step 2 --- Producer on CM33\_0 (1\,kHz ADC Task)}

\begin{lstlisting}[style=cstyle]
void vADCTask(void *pv)
{
    TickType_t xLastWake = xTaskGetTickCount();
    uint32_t sample_count = 0;
    for (;;) {
        vTaskDelayUntil(&xLastWake, pdMS_TO_TICKS(1));  /* 1 kHz */

        uint32_t next = (g_ring.write_idx + 1) % RING_BUF_SIZE;
        if (next == g_ring.read_idx) {
            g_ring.missed++;
            continue;   /* buffer full -- drop sample */
        }

        g_ring.data[g_ring.write_idx] = (int16_t)(sample_count % 4096);
        sample_count++;
        __DMB();                        /* store before index update */
        g_ring.write_idx = next;
        __DMB();                        /* index update before MU notify */

        if ((g_ring.write_idx & 0xF) == 0)
            MU_TrySendMsg(MUA, 0, g_ring.write_idx);
    }
}
\end{lstlisting}

\subsection{Step 3 --- Consumer on CM33\_1 (FIR Filter Task)}

\begin{lstlisting}[style=cstyle]
static int16_t fir_h[8] = {1, 2, 3, 4, 4, 3, 2, 1};
static int32_t fir_state[8] = {0};

void MUA_IRQHandler(void)
{
    MU_ReceiveMsg(MUB, 0);   /* consume notification */
    MU_ClearStatusFlags(MUB, kMU_Rx0FullFlag);
    xTaskNotifyGiveFromISR(xFIRHandle, NULL);
}

void vFIRTask(void *pv)
{
    uint32_t total = 0;
    for (;;) {
        ulTaskNotifyTake(pdTRUE, portMAX_DELAY);
        __DMB();

        while (g_ring.read_idx != g_ring.write_idx) {
            int16_t sample = g_ring.data[g_ring.read_idx];
            g_ring.read_idx = (g_ring.read_idx + 1) % RING_BUF_SIZE;

            memmove(&fir_state[1], &fir_state[0], 7 * sizeof(int32_t));
            fir_state[0] = sample;
            int32_t acc = 0;
            for (int k = 0; k < 8; k++) acc += fir_state[k] * fir_h[k];
            total++;
        }

        if ((total % 1000) == 0)
            uart_printf("[CM1] processed %lu samples  missed=%lu\r\n",
                        total, g_ring.missed);
    }
}
\end{lstlisting}

\begin{checkpointbox}[~C]
  UART from CM33\_1 showing processed sample count and zero missed samples at 1\,kHz.
\end{checkpointbox}

\begin{questionbox}[~C1]
  Remove both \api{\_\_DMB()} calls on the producer side, rebuild, and run. Do you
  observe incorrect FIR outputs or missed samples? Re-add the barriers before
  continuing.
\end{questionbox}

\begin{questionbox}[~C2]
  The ring buffer uses \texttt{volatile} on the indices. Is \texttt{volatile}
  sufficient to guarantee correctness on a multi-core system? What does
  \texttt{volatile} guarantee vs.\ what \api{\_\_DMB()} guarantees?
\end{questionbox}

% ── Part D ─────────────────────────────────────────────────────────────────────
\section{Part D --- Comparison: MU Direct vs.\ Ring Buffer}

Run both configurations with a 1\,kHz producer for 10 seconds. Fill in the table:

\renewcommand{\arraystretch}{1.4}
\begin{tabular}{@{} p{5.5cm} p{3.5cm} p{3.5cm} @{}}
  \toprule
  \textbf{Metric} & \textbf{MU direct (Part B)} & \textbf{Ring buffer (Part C)} \\
  \midrule
  Max throughput (samples/s) & & \\
  Latency per sample (cycles) & & \\
  Missed samples at 1\,kHz    & & \\
  CM33\_0 CPU load (\%)       & & \\
  \bottomrule
\end{tabular}
\renewcommand{\arraystretch}{1}

\medskip
Measure CM33\_0 CPU load using \api{vTaskGetRunTimeStats} with
\texttt{configGENERATE\_RUN\_TIME\_STATS\,=\,1} and DWT as the timer source.

\begin{checkpointbox}[~D]
  Completed comparison table with measured values.
\end{checkpointbox}

% ── Deliverables ───────────────────────────────────────────────────────────────
\section{Deliverables}

\renewcommand{\arraystretch}{1.3}
\begin{tabular}{@{} c p{5cm} p{6.5cm} @{}}
  \toprule
  \textbf{\#} & \textbf{Item} & \textbf{Details} \\
  \midrule
  1 & Part A UART & Both core startup messages \\
  2 & Part B2 UART & MU latency in cycles + nanoseconds \\
  3 & Part C UART & 10 seconds of CM33\_1 processing, zero missed \\
  4 & Linker script excerpt & \texttt{.sram4} section placement \\
  5 & Part D comparison table & Filled with measured values \\
  6 & Written answers & Questions A1, B1, C1, C2 \\
  7 & \file{shared\_mem.h} & Ring buffer struct listing \\
  \bottomrule
\end{tabular}
\renewcommand{\arraystretch}{1}

% ── Key Concepts ───────────────────────────────────────────────────────────────
\section{Key Concepts Demonstrated}

\renewcommand{\arraystretch}{1.3}
\begin{tabular}{@{} p{6.5cm} p{6cm} @{}}
  \toprule
  \textbf{Concept} & \textbf{Where you saw it} \\
  \midrule
  AMP dual-core boot via SYSCON & Part A \\
  MU direct 32-bit messaging & Part B \\
  IPC latency measurement with DWT & Part B \\
  Shared SRAM ring buffer & Part C \\
  DMB memory barriers & Part C \\
  Throughput vs.\ latency trade-off & Part D \\
  \bottomrule
\end{tabular}
\renewcommand{\arraystretch}{1}

% ── Reference ─────────────────────────────────────────────────────────────────
\section{Reference}

\renewcommand{\arraystretch}{1.3}
\begin{tabular}{@{} p{6cm} p{6.5cm} @{}}
  \toprule
  \textbf{Resource} & \textbf{Relevant sections} \\
  \midrule
  NXP MCXN236 Reference Manual & Chapter MU (Messaging Unit), Chapter SYSCON \\
  NXP MCUXpresso SDK & \file{fsl\_mu.h}, multicore examples \\
  ARM Cortex-M33 TRM & \S A3.4 Memory ordering \\
  Barry, \textit{Mastering the FreeRTOS Kernel} &
    \api{xTaskNotifyGiveFromISR}, \api{vTaskGetRunTimeStats} \\
  \bottomrule
\end{tabular}
\renewcommand{\arraystretch}{1}

\end{document}
